\section*{ID6183: Add a Toroidal (Orbital) Gaussian Envelope: G\_ell(r,θ)}
\paragraph{Metadata.}
PiID: 9009019211697 | RTEID: RTE:P24009.9621:T1.8619 | UUID: 36770c13-ebd7-58d9-b423-6b5610c43ca6

\paragraph{Canonical Statement.}
\# Add a toroidal (orbital) Gaussian envelope A toroidal / ring / vortex Gaussian in the **relative** coordinate is naturally written in polar form when the relative coordinate is a 2D vector \textbackslash{}(\textbackslash{}mathbf r=(r,\textbackslash{}theta)\textbackslash{}). A common form (Laguerre–Gaussian–like) is \textbackslash{}[ G\_\textbackslash{}ell(r,\textbackslash{}theta)=\textbackslash{}mathcal N\textbackslash{}, r\textasciicircum{}\{|\textbackslash{}ell|\}\textbackslash{}, e\textasciicircum{}\{-r\textasciicircum{}2/(2\textbackslash{}sigma\_r\textasciicircum{}2)\} e\textasciicircum{}\{i\textbackslash{}ell\textbackslash{}theta\}, \textbackslash{}] where \textbackslash{}(\textbackslash{}ell\textbackslash{}in\textbackslash{}mathbb Z\textbackslash{}) is the orbital angular momentum (OAM) index, \textbackslash{}(\textbackslash{}sigma\_r\textbackslash{}) the ring width, and \textbackslash{}(\textbackslash{}mathcal N\textbackslash{}) a normalization. For a 1D model you can use the radial envelope with the absolute value: \textbackslash{}[ G\_\textbackslash{}ell(r)=\textbackslash{}mathcal N\textbackslash{},|r|\textasciicircum{}\{|\textbackslash{}ell|\} e\textasciicircum{}\{-r\textasciicircum{}2/(2\textbackslash{}sigma\_r\textasciicircum{}2)\} . \textbackslash{}]

\paragraph{Canonical Equation.}
\[
G_\ell(r,\theta)=\mathcal N\, r^{|\ell|}\, e^{-r^2/(2\sigma_r^2)} e^{i\ell\theta},
\]

\paragraph{Dictionary.}
Add a Toroidal (Orbital) Gaussian Envelope: G\_ell(r,θ) — G\_ℓ(r,θ)=mathcal N r\textasciicircum{}|ℓ| e\textasciicircum{}-r\textasciicircum{}2/(2σ\_r\textasciicircum{}2) e\textasciicircum{}iℓθ

\paragraph{Encyclopedia.}
Add a Toroidal (Orbital) Gaussian Envelope: G\_ell(r,θ) is an exponential / gaussian profile, written G\_ℓ(r,θ)=mathcal N r\textasciicircum{}|ℓ| e\textasciicircum{}-r\textasciicircum{}2/(2σ\_r\textasciicircum{}2) e\textasciicircum{}iℓθ. It describes an exponential or Gaussian profile characterizing localized or decaying behaviour. It is built from ell, θ, N, r, defining G\_ℓ(r,θ). Within the Trawin Zero Point registry it serves as one element of the framework's mathematical narrative.
